\documentclass[11pt,a4paper]{article}
\usepackage[margin=25mm]{geometry}
\usepackage[T1]{fontenc}
\usepackage[utf8]{inputenc}
\usepackage{lmodern,microtype}
\usepackage{amsmath,amssymb,amsthm,mathtools}
\usepackage{booktabs,longtable,array,multicol,enumitem}
\usepackage[hidelinks]{hyperref}
\allowdisplaybreaks
\newtheorem{theorem}{Theorem}[section]
\newtheorem{lemma}[theorem]{Lemma}
\newtheorem{proposition}[theorem]{Proposition}
\newtheorem{corollary}[theorem]{Corollary}
\theoremstyle{remark}
\newtheorem{remark}[theorem]{Remark}
\newcommand{\F}{\mathbb F_2}
\newcommand{\K}{\mathcal K}
\newcommand{\E}{\mathcal E}
\newcommand{\Forb}{\operatorname{Forb}}
\newcommand{\sgn}{\operatorname{sgn}}
\newcommand{\supp}{\operatorname{supp}}
\newcommand{\bit}[1]{\mathtt{#1}}
\title{Negative amalgamation results for $(n+1)$-constrained hypertournaments\\\large Formalisation-preflight audit and Lean-oriented technical note}
\author{ChatGPT (GPT-6)\\\small AI-generated working note for Mat\v ej Kone\v cn\'y\\\small Not independently refereed or formally verified}
\date{6 November 2026 --- formalisation preflight v4}
\begin{document}
\maketitle
\begin{abstract}
We give a detailed, self-contained account of the \emph{negative}
$(n+1)$-constrained amalgamation classification for finite
$n$-hypertournaments, $n\ge4$. No positive amalgamation theorem is used.
The proof has explicit non-amalgamable seeds, a one-coordinate cone,
two type-shifting suspensions, and two uniform obstruction families.
The two suspensions now share one general reflection lemma, and the
probabilistic construction includes an explicit finite parameter bound.
Two finite seed calculations are accompanied by exhaustive, independent
checking scripts; the other arguments are mathematical proofs rather
than computer assertions.
This is an AI-written mathematical working note, checked in detail but
not independently refereed or Lean-verified. Its intended uses are
human mathematical review and subsequent formalisation. The separate
five-constrained ternary census and its electronic certificates are
not reproduced here.
\end{abstract}
\tableofcontents

\section{Conventions and negative classification statements}\label{sec:statement}
An \emph{orientation} of an $n$-set is an ordering of its vertices
up to even permutation. An \emph{$n$-hypertournament} $B$ assigns one
of its two orientations to every $n$-element subset of a finite vertex
set $V(B)$. Embeddings are injections preserving these orientations.
The induced substructure on $U\subseteq V(B)$ is given by restriction.

\paragraph{Amalgamation conventions.}
An \emph{ordinary amalgam} of embeddings $A\hookrightarrow B_\ell$
and $A\hookrightarrow B_r$ consists of a hypertournament $C$ and
embeddings $f_\ell:B_\ell\hookrightarrow C$,
$f_r:B_r\hookrightarrow C$ agreeing on $A$.
The property \emph{AP} means that every such diagram has an amalgam.
The amalgam is \emph{strong} if the two images intersect exactly in
(the image of) $A$. For a \emph{one-point diagram}
$B_\ell=A\cup\{\ell\}$ and $B_r=A\cup\{r\}$, the two new
vertices of any ordinary amalgam are either identified or remain
distinct. In the second case, restriction to the union of the images
produces a strong amalgam on $A\cup\{\ell,r\}$; it is therefore
enough to exclude these two possibilities. Our sign computations
always use one common auxiliary order on that union.

\paragraph{Normalized facet packets.}
Fix an auxiliary linear order, which is \emph{not} part of the
language. Write $s_B(F)\in\{+1,-1\}$ for the sign of the chosen
orientation of $F$, relative to the increasing enumeration. If
$X=\{x_0<\cdots<x_n\}$, define
\[
 t_i(X)=(-1)^i s_B(X\setminus\{x_i\}),\qquad
 d_n(X)=\min\bigl(\#\{i:t_i=+1\},\#\{i:t_i=-1\}\bigr).
\]
Changing the auxiliary order permutes the normalized signs and may
reverse all of them simultaneously. Consequently $d_n$ is intrinsic.
Conversely, two $(n+1)$-vertex hypertournaments have the same $d_n$
if and only if they are isomorphic: permutations of the facets act on
normalized packets by signed permutations, and an odd permutation
interchanging two equally signed coordinates can adjust the parity
of the relabelling (there are always at least two such coordinates).
For a precise coordinate rule, if a relabelling permutes the
$(n+1)$ vertices by $\pi$, then its action on the normalised packet is
\[
  t'_{\pi(i)}=\operatorname{sgn}(\pi)t_i.
\]
Indeed the parity of the permutation induced on the facet omitting
vertex $i$ is $\operatorname{sgn}(\pi)(-1)^{i+\pi(i)}$.
This identity proves intrinsicness without choosing a preferred
order. For completeness of the invariant, first match the sets of
positive coordinates (possibly after reversing all signs), then
adjust the parity of the matching permutation by exchanging two
coordinates of the same sign; such coordinates exist when $n+1\ge3$.
The isomorphism types are thus $T_{n,d}$ for
$0\le d\le q=\lfloor(n+1)/2\rfloor$. We write
$H_{n+1}=T_{n,0}$ and $C_{n+1}=T_{n,q}$.

For a set $S$ of local types, let $\K^{(n)}_S$ consist of finite
$n$-hypertournaments $B$ such that every $(n+1)$-subset of $V(B)$
has one of the types in $S$. A packet is \emph{central} when it has
type $T_{n,q}$ and \emph{homogeneous} when it has type $T_{n,0}$.

For context only, the finite $n$-uniform Ramsey theorem shows that $\K^{(n)}_S$ is unbounded if and
only if $T_{n,q}\in S$: after colouring the increasing $n$-sets by their
orientation, any sufficiently large structure contains an
$(n+1)$-set of constant colour and hence of central type; conversely,
the increasing orientation on arbitrarily many vertices realizes only
central types. \emph{This Ramsey fact is not used in the negative
implications: throughout the proof, the hypothesis $T_{n,q}\in S$
is taken explicitly.}

\begin{theorem}[Necessary conditions for AP in the $(n+1)$-constrained case]
\label{thm:classification}
Let $n\ge4$, set $q=\lfloor(n+1)/2\rfloor$, and assume that
$T_{n,q}\in S$ (equivalently, $\K^{(n)}_S$ is unbounded).
\begin{enumerate}[label=\textup{(\roman*)},leftmargin=15mm]
\item If $n$ is even and $T_{n,q-1}\notin S$, then
$\K^{(n)}_S$ fails ordinary AP.
\item If $n\ge5$ is odd, $T_{n,q-1}\notin S$, and
\[
 S\ne E_n:=\{T_{n,d}: d\equiv q\pmod2\},
\]
then $\K^{(n)}_S$ fails ordinary AP.
\end{enumerate}
In each negative case a one-point-versus-one-point
non-amalgamation diagram exists, with the possibility of identifying
the new vertices explicitly excluded.
\end{theorem}

\begin{corollary}[Exportable negative-classification interface]
\label{cor:export}
Let $n\ge4$, let $S\subseteq\{T_{n,0},\ldots,T_{n,q}\}$
contain the central type $T_{n,q}$, and suppose that
$\K_S^{(n)}$ has ordinary AP. Then
\[
 \boxed{T_{n,q-1}\in S\quad\text{or}\quad
      (n\text{ odd and }S=E_n).}
\]
\end{corollary}
\begin{proof}
If $T_{n,q-1}\in S$, there is nothing to prove. Otherwise
Theorem~\ref{thm:classification}(i) rules out even $n$. Thus $n$
is odd (and $n\ge5$), and part (ii) of that theorem forces
$S=E_n$.
\end{proof}

The explicit inclusion $S\subseteq\{T_{n,0},\ldots,T_{n,q}\}$
is important for machine formalisation. An arbitrary subset of
$\mathbb N$ may contain indices which are \emph{not} local types;
adding such indices changes neither the class nor AP but would
invalidate the literal equality $S=E_n$. In Lean we either use a
bounded index type or carry this condition as a separate hypothesis.

The families outside the conclusion are listed solely as \emph{exceptions
to the negative theorem}; their positive AP/SAP proofs belong to the
companion manuscript~\cite{catalogue} and are not reproduced here.
For arities $n=2,3$ there is no negative case in the unbounded
$(n+1)$-constrained catalogue.

For clarity we keep two distinct conventions for Boolean data. In the
arity-four and arity-five seeds, $x=1$ means \emph{positive} sign;
in the parity constructions below, $e=0$ means positive sign.
The distinction is explicitly restated when needed, and all arithmetic
on $\mathbb F_2$ is modulo $2$.

For odd $n$ it is convenient to encode an orientation by a bit
$e(F)\in\F$, with $e(F)=0$ denoting the positive increasing sign, and
$e(F)=1$ the negative sign.  The convention is used consistently within
the parity sections; the sign-based seed arguments will explicitly state
their own Boolean convention.  The unsigned coboundary over $\F$ is
\[
 (\delta f)(Y)=\sum_{x\in Y}f(Y\setminus\{x\})\pmod2,
\qquad \delta^2=0.
\]
For odd $n$, the switching class $E_n$ is exactly the class in which
$\delta e=0$ on every $(n+1)$-set.  Indeed the parity of a raw
$(n+1)$-facet packet equals $d_n+(n+1)/2$ modulo two.  We will use the following elementary exactness statement. It is
included to avoid any dependence on a topological homology theorem.

\begin{lemma}[A contracting cochain on a simplex]\label{lem:exactness}
Let $V$ be finite, $r\ge2$, and let
$g:\binom V r\to\mathbb F_2$ satisfy $\delta g=0$ on each
$(r+1)$-set. Fix $a\in V$ and define
$u:\binom V{r-1}\to\mathbb F_2$ by
\[
 u(T)=\begin{cases}g(\{a\}\cup T),&a\notin T,\\0,&a\in T.\end{cases}
\]
Then $g=\delta u$. In addition, $\delta^2=0$ in every arity.
\end{lemma}
\begin{proof}
Each term in $\delta^2$ occurs exactly twice, so it vanishes in
$\mathbb F_2$. If $a\in F\in\binom Vr$, only the deletion of $a$
contributes to $\delta u(F)$, giving $g(F)$. If $a\notin F$,
the equation $\delta g(\{a\}\cup F)=0$ gives
$g(F)=\sum_{v\in F}g(\{a\}\cup(F\setminus\{v\}))=\delta u(F)$.
\end{proof}

For odd $n$, $\delta e=0$ on all $(n+1)$-sets therefore gives
$e=\delta z$ for an $(n-1)$-set potential whenever the ambient
vertex set is nonempty. (When fewer than $n$ vertices are present,
there are no $n$-set orientations to reconstruct.)

For integer-valued functions on fixed-size subsets we use the
\emph{alternating} boundary operator
\[
 (\partial f)(\{v_0<\cdots<v_r\})
 =\sum_{i=0}^r(-1)^if(\{v_0,\ldots,v_r\}\setminus\{v_i\}).
\]
Here $\partial^2=0$ by pairwise cancellation, just as for $\delta$.

\section{Two elementary arity-four obstructions}\label{sec:arity4}
We use $A=\{0,\ldots,m-1\}$, with a new vertex $\ell<A<r$.
For $n=4$, put $x_{ij}=s(\{\ell,i,j,r\})\in\{\pm1\}$ and
write $\alpha(ijk)=s(\ell ijk)$, $\beta(ijk)=s(ijkr)$.
The mixed five-set $\{\ell,i,j,k,r\}$ has normalized facet signs
\begin{equation}\label{eq:4-packet}
 \bigl(\beta(ijk),-x_{jk},x_{ik},-x_{ij},\alpha(ijk)\bigr).
\end{equation}

\subsection{The central-only class $\K^{(4)}_{\{T_{4,2}\}}$}
\begin{lemma}\label{lem:4-central}
The class of finite $4$-hypertournaments in which every five-set has
central type $T_{4,2}$ has no ordinary AP.
\end{lemma}
\begin{proof}
Set $A=\{0,\ldots,6\}$.  For an ordered four-set, list below those
with sign $+1$; all unlisted signs are $-1$:
\begin{align*}
\supp(s|_A)&=\{0123,0125,0234,0236,0256,1235,1245\},\\
\supp(\alpha)&=\{013,014,015,016,024,025,026,035,045,046,056,156,\\[-2pt]
 &\hspace{43mm}245,256,356,456\},\\
\supp(\beta)&=\{014,015,016,025,046,056,146,156,246,256,346,356,456\}.
\end{align*}
The given base and both one-point extensions are central-only: on
every one of their five-subsets, the normalized vector has two or three
positive signs.  This is also an immediate check using the three
finite lists above.

Suppose a strong amalgam exists, and write $x_{ij}\in\{0,1\}$ for the
truth value of ``the mixed four-set has sign $+1$''.  On each triple
$Q=ijk$ with $\alpha(Q)=\beta(Q)=-1$, the central condition
\eqref{eq:4-packet} is equivalent to
\[
 x_{ij}\Rightarrow x_{ik},\qquad
 x_{jk}\Rightarrow x_{ik},\qquad
 \neg(x_{ij}\wedge x_{jk}). \tag{$-$}
\]
When both endpoint signs are $+1$, it is equivalent to
\[
 x_{ik}\Rightarrow x_{ij},\qquad
 x_{ik}\Rightarrow x_{jk},\qquad
 x_{ij}\vee x_{jk}. \tag{$+$}
\]
For the seven triples
\[
 012,\quad014,\quad023,\quad036,\quad046,\quad145,\quad456,
\]
the common endpoint sign is respectively
$-,+,-,-,+,-,+$.  Using only these implications gives
\[
 \neg x_{01}\Rightarrow\neg x_{04}\Rightarrow x_{46}
 \Rightarrow x_{45}\Rightarrow\neg x_{14}\Rightarrow x_{01}
\]
and
\[
\begin{split}
 x_{01}&\Rightarrow x_{02}\Rightarrow x_{03}\Rightarrow x_{06}
 \Rightarrow x_{04}\Rightarrow x_{14}\Rightarrow\neg x_{45}\\
 &\Rightarrow\neg x_{46}\Rightarrow\neg x_{06}
 \Rightarrow\neg x_{03}\Rightarrow\neg x_{02}
 \Rightarrow\neg x_{01}.
\end{split}
\]
Both truth values of $x_{01}$ lead to a contradiction.  An ordinary
amalgam cannot identify $\ell$ and $r$, because such an identification
would require $\alpha=-\beta$, whereas their displayed links do not
satisfy this equality.  Hence ordinary AP fails.
\end{proof}

\subsection{The class omitting $T_{4,1}$}
\begin{lemma}\label{lem:4-forb}
The class $\K^{(4)}_{\{T_{4,0},T_{4,2}\}}=\Forb(T_{4,1})$
does not have ordinary AP.
\end{lemma}
\begin{proof}
Use $A=\{0,1,2,3,4,5\}$, with just the base four-set $1245$
positive.  Put
\[
 \supp(\alpha)=\{035,135,235,345\},\qquad
 \supp(\gamma)=\{012,034,035,045,134,245,345\},\quad
 \gamma=-\beta.
\]
Every base five-set is central; among those using $\ell$, precisely
$\ell0345$, $\ell1345$ and $\ell2345$ are homogeneous and all others
central; every five-set using $r$ is central.  Thus the sides omit
$T_{4,1}$.

As above, set $v_{ij}=1$ when $x_{ij}=+1$ and $v_{ij}=0$ otherwise.
On a triple $ijk$ with $i<j<k$, the condition that the mixed packet
\eqref{eq:4-packet} omit $T_{4,1}$ is a constraint on the three bits
$(v_{ij},v_{ik},v_{jk})$.  It falls into just three types:
\[
\begin{array}{c|l}
 \text{type}&\text{forbidden triples of bits}\\\hline
 \mathrm T&010,101\\
 \mathrm C_-&001,100,111\\
 \mathrm C_+&000,011,110
\end{array}
\]
Here $\mathrm T$ means $\alpha=\gamma$, and $\mathrm C_\pm$
means $\alpha=-\gamma=\pm1$.  From the displayed links the
restrictions are
\[
\begin{array}{r|l}
 \mathrm T&013,014,015,023,024,025,035,123,124,125,145,234,345\\
 \mathrm C_-&012,034,045,134,245\\
 \mathrm C_+&135,235.
\end{array}\tag{*}
\]

We give a finite \emph{implication proof}, not an appeal to a search
certificate.  First restrict $v$ to $\{1,2,3,4,5\}$.  The ten
triples inside this set, using the rule $(*)$, have exactly the
seventeen possible assignments displayed in Table~\ref{tab:n4paths}.
For a direct check of exhaustiveness, Table~\ref{tab:n4extend}
shows all sixteen assignments on $\{1,2,3,4\}$ that survive its
four triples, and all their permitted choices of the four incident
edges to vertex $5$; these are obtained by applying $(*)$ to its
six new triples.  Thus no other assignment is missing.

Fix one of those seventeen restrictions and let $y_i=v_{0i}$.
For $i<j$ in $\{1,\ldots,5\}$, the mixed triple $0ij$ gives the
following $2$-variable restriction (where a pair such as $01$
denotes $(y_i,y_j)=(0,1)$):
\[
\begin{array}{c|cc}
 &v_{ij}=0&v_{ij}=1\\\hline
 ij\in\{12,34,45\}&(10)\text{ forbidden}&y_i\ne y_j\\
 \text{other }ij &(01)\text{ forbidden}&(10)\text{ forbidden}.
\end{array}\tag{**}
\]
In the exceptional cell $v_{ij}=1$, $ij\in\{12,34,45\}$,
$y_i\ne y_j$ is a \emph{required} relation; every other cell
lists a forbidden pair. Each forbidden pair gives an elementary
implication from either literal to the negation of the other.  Table~\ref{tab:n4paths}
exhibits, in every one of the seventeen cases, both an implication
path $\neg y_i\Rightarrow y_i$ and a path
$y_i\Rightarrow\neg y_i$, using only $(**)$.  Every path step is
between the two literals indexed by the pair $ij$ that supplies its
clause; hence each step can be checked directly from the table $(**)$.
This is a contradiction independent of any remaining bits.

Identification of $\ell,r$ would force $\alpha=\gamma$ on every
triple, contrary to the explicit lists.  So the diagram has no
ordinary amalgam.
\end{proof}

\section{Arity five: a central cone and a defect-graph obstruction}
\label{sec:arity5}
The class $\K^{(5)}_{\{T_{5,0},T_{5,3}\}}$ is a useful extra
seed.  We give a mathematical proof of its failure of AP; it is not
necessary to appeal to its earlier exhaustive-computation certificate.
The third negative seed in arity five, $\K^{(5)}_{\{T_{5,3}\}}$,
will follow from a one-coordinate cone, and the largest negative seed,
$\K^{(5)}_{\{T_{5,0},T_{5,1},T_{5,3}\}}$, from the uniform
obstruction in Section~\ref{sec:one-four}.

\begin{lemma}[The zero-or-two principle]\label{lem:zero-two}
Let $n\equiv1\pmod4$ and $n\ge5$, and suppose every $(n+1)$-set
in an $n$-hypertournament is either central to the switching parity
or homogeneous; more precisely, its type belongs to
$\{T_{n,0},T_{n,1},T_{n,3},\ldots,T_{n,(n+1)/2}\}$.
Then every $(n+2)$-vertex subset contains either zero or exactly two
homogeneous $(n+1)$-subsets.
\end{lemma}
\begin{proof}
Put $N=n+1\equiv2\pmod4$.  In the bit convention of
Section~\ref{sec:statement}, a nonhomogeneous permitted $N$-packet
has even raw parity, while a homogeneous packet is alternating and
has odd raw parity.  Thus its indicator is
$D(X)=\sum_{F\in\binom Xn}e(F)\pmod2$.
On $N+1$ vertices the number of homogeneous $N$-subsets is even:
every $n$-subset contributes to two $N$-subsets, so
$\delta^2e=0$.

It is at most two.  Write $Z=\{z_0<\cdots<z_N\}$ and suppose
$Z\setminus\{z_i\}$ is homogeneous, whose raw packet has $e$-bits $(\lambda_i+j\bmod2)_j$,
with phase $\lambda_i\in\F$. Comparing two such homogeneous subsets,
indexed $i<j$, on their common $n$-facet gives
\[
 \lambda_i+\lambda_j=i+j+1\pmod2.
\]
Equivalently, $\mu_i:=\lambda_i+i$ satisfies
$\mu_i+\mu_j=1$ for every pair of homogeneous subsets.
Three such subsets yield three mutually inconsistent equations over
$\F$.  Therefore there are at most two, and parity finishes.
\end{proof}

\subsection{The class $\{T_{5,0},T_{5,3}\}$}
\begin{lemma}\label{lem:5-two}
The $6$-constrained class
$\K^{(5)}_{\{T_{5,0},T_{5,3}\}}$ fails ordinary AP.
\end{lemma}
\begin{proof}
Let $A=\{0,1,2,3,4,5,6\}$ and put $\ell<A<r$.
Encode the positive signs on base five-sets and on the two four-set
links by the following lists (all unlisted signs are negative):
\begin{align*}
\supp(s|_A)&=\{01235,01246,01356,02345,02456,12346,13456\},\\
\supp(\alpha)&=\{0123,0124,0125,0126,0135,0145,0146,0234,0245,\\[-2pt]
 &\hspace{35mm}0345,0346,1234,1345\},\\
\supp(\beta)&=\{0124,0134,0145,0156,0234,0236,0256,0346,0356,\\[-2pt]
 &\hspace{20mm}1234,1235,1236,1245,1345,1356,1456,2346,2356\}.
\end{align*}
Here $\alpha(Q)=s(\{\ell\}\cup Q)$ and
$\beta(Q)=s(Q\cup\{r\})$.  A direct facet check shows that all
seven six-sets of $A$ have type $T_{5,3}$, all twenty-one six-sets
of $A\cup\{\ell\}$ containing $\ell$ have type $T_{5,3}$, and
in $A\cup\{r\}$ precisely three six-sets are homogeneous, namely
\[
 (A\setminus\{2,3\})\cup\{r\},\quad
 (A\setminus\{2,6\})\cup\{r\},\quad
 (A\setminus\{3,6\})\cup\{r\};
\]
the remaining six-sets on the right are central.  Thus both sides
lie in the asserted class.

Suppose there is a strong amalgam on
$V=A\cup\{\ell,r\}$.  For a triple $T\subseteq A$ write
$y_T=1$ when the mixed six-set $V\setminus T$ is homogeneous, and
put $y_T=0$ otherwise.  The two fixed links can occur at opposite
ends of an alternating packet only if
$\alpha(A\setminus T)\ne\beta(A\setminus T)$.  Therefore $y_T$
may be nonzero only for
\begin{equation}\label{eq:eligible5}
\begin{split}
\mathcal U=\{&014,015,023,024,036,045,046,124,126,134,136,145,\\
 &234,235,246,256,345,346,456\}.
\end{split}
\end{equation}
By Lemma~\ref{lem:zero-two}, for every pair
$ij\in\binom A2$, the number of homogeneous six-sets in
$V\setminus\{i,j\}$ is $0$ or $2$.  The right-side homogeneous
six-sets contribute once exactly for pairs $23,26,36$.
Thus, writing $h_{ij}=\mathbf1_{ij\in\{23,26,36\}}$, we have
\begin{equation}\label{eq:pair-degree5}
 h_{ij}+\sum_{T\in\mathcal U:\ ij\subseteq T}y_T\in\{0,2\}
 \qquad (ij\in\tbinom A2).
\end{equation}
The sum in~\eqref{eq:pair-degree5} is an \emph{integer} sum,
not an $\mathbb F_2$ sum. We now solve these small incidence
constraints explicitly,
so that no computer assertion enters the proof.

First reduce~\eqref{eq:pair-degree5} modulo two.  Take as four free
bits
\[
 a=y_{014},\qquad b=y_{023},\qquad c=y_{124},\qquad d=y_{234}.
\]
The pair equations, in the order
\[
 01,02,03,05,04,12,15,14,13,23,24,25,35,34,45,
\]
uniquely determine the other fifteen bits: at every stage the
indicated pair equation has exactly one as-yet-unknown variable.
Substituting into the remaining pair equations gives identities.
Equivalently, with symmetric difference of sets, the resulting
support is
\begin{equation}\label{eq:four-free5}
 \mathcal S(a,b,c,d)
 =S_*\ \triangle\ aU_a\ \triangle\ bU_b\ \triangle\ cU_c\ \triangle\ dU_d,
\end{equation}
where
\begin{align*}
 S_*&=\{235,256,345,346,456\},\\
 U_a&=\{014,015,045,145\},\\
 U_b&=\{023,024,036,046,235,246,256,345,346,456\},\\
 U_c&=\{124,126,134,136,246,346\},\\
 U_d&=\{234,235,246,256,345,456\}.
\end{align*}
There are only sixteen choices of $(a,b,c,d)$.  For eleven of them
some pair has incidence count four in~\eqref{eq:pair-degree5},
contradicting the upper bound two.  Table~\ref{tab:n5supports}
records an offending pair for each of these eleven choices.
The other five possibilities are, writing each triple by its three
digits,
\[
\begin{array}{c|l}
 \mathrm A&124,126,134,136,234\\
 \mathrm B&023,024,036,046,246\\
 \mathrm C&234,246,346\\
 \mathrm D&014,015,045,145,234,246,346\\
 \mathrm E&235,256,345,346,456.
\end{array}\tag{$\dagger$}
\]

It remains to rule out these five cases.  Write
$x_{ijk}\in\F$ for the raw orientation bit of the mixed five-set
$\{\ell,i,j,k,r\}$, with $x=1$ for positive sign, and let
$\alpha(Q),\beta(Q)$ also denote $0$--$1$ positive-sign bits.
For $Q=\{a<b<c<d\}$ the six normalized positive-sign bits are
\begin{equation}\label{eq:mixed5packet}
 \bigl(\beta(Q),1+x_{bcd},x_{acd},1+x_{abd},x_{abc},1+\alpha(Q)\bigr).
\end{equation}
When $V\setminus(A\setminus Q)$ is not homogeneous, this packet must
have \emph{exactly three} entries equal to one.  When it is
homogeneous, its entries are all equal, and the first coordinate
fixes every mixed facet bit.  The following deductions use only this
observation and the lists defining the sides:
\begin{enumerate}[label=\textup{(\Alph*)},leftmargin=15mm]
\item In case A, homogeneity fixes
$x_{035}=x_{015}=1$, $x_{036}=x_{026}=0$.
Centrality on $Q=0135$ forces $x_{013}=1$; centrality on $0236$
forces $x_{023}=0$.  But on $0123$, the normalized packet has
at most two positive entries, irrespective of its other bits.
\item In case B, homogeneity fixes
$x_{123}=1$, $x_{013}=0$, $x_{156}=x_{015}=1$, $x_{136}=0$.
Centrality on $0123$ forces $x_{023}=1$;
on $0156$ it forces $x_{016}=1$;
on $0236$ it forces $x_{036}=0$.
The packet on $0136$ then has exactly two positive entries,
a contradiction.
\item In case C, homogeneity fixes
$x_{013}=x_{012}=0$, $x_{056}=1$, $x_{035}=0$.
Centrality on $0123$ forces $x_{023}=1$; centrality on $0356$
forces $x_{036}=1$.  The packet on $0236$ has at least four
positive entries.
\item In case D, homogeneity already forces
$(x_{123},x_{023},x_{013},x_{012})=(1,1,0,0)$.
The packet on $0123$ has only two positive entries.
\item In case E, homogeneity forces
$x_{015}=x_{013}=1$, $x_{026}=x_{023}=0$.
Centrality on $0135$ gives $x_{035}=1$, while centrality on $0236$
gives $x_{036}=0$.  The packet on $0356$ then has at least four
positive entries.
\end{enumerate}
All five possible homogeneous-defect configurations are impossible.
Thus no strong amalgam exists.  The vertices $\ell,r$ cannot be
identified: their four-set links $\alpha,\beta$ disagree, and
moving the new vertex across the four retained base vertices
changes no sign.  Hence the obstruction is to ordinary AP.
\end{proof}

\subsection{A cone from even central arity to odd central arity}
\begin{proposition}[Central-only cone]\label{prop:cone}
Let $n\ge4$ be even.  If $\K^{(n)}_{\{T_{n,n/2}\}}$ has a
one-point diagram without ordinary amalgam, so does
$\K^{(n+1)}_{\{T_{n+1,n/2+1}\}}$.
\end{proposition}
\begin{proof}
Write $s$ for the original signed orientation and
$\partial s(Y)=\sum_{i=0}^{n}(-1)^is(Y\setminus\{y_i\})$
for the alternating integer coboundary.  For every old central
$(n+1)$-set $Y$, its coboundary is $\pm1$.
Adjoin one common vertex $p$ before all the old vertices and put
\[
 s'(Z)=\begin{cases}
 s(Z\setminus\{p\}),&p\in Z,\\
 \partial s(Z),&p\notin Z,
 \end{cases}\qquad |Z|=n+1.
\]
On an old $(n+2)$-set, $\partial s'=\partial^2s=0$.
On $\{p\}\cup Y$, the first boundary term is $\partial s(Y)$,
and the other terms add up to $-\partial s(Y)$.
So every new $(n+2)$-set is central.
Conversely, in any proposed amalgam, orient each old $n$-set
$F$ by the sign of $\{p\}\cup F$.  Centrality on
$\{p\}\cup Y$ forces the old coboundary to be the orientation
of $Y$, hence to have absolute value one.  This recovers an old
central-only amalgam, including any attempted identification of the
new vertices.  Contradiction.
\end{proof}

In particular, Lemma~\ref{lem:4-central} and
Proposition~\ref{prop:cone} give a purely mathematical negative
seed for $\K^{(5)}_{\{T_{5,3}\}}$.  All three negative
arity-five seeds now have mathematical proofs once the uniform
$1\pmod4$ construction below is included.


\section{A common reflection lemma for two-coordinate lifts}
\label{sec:reflection}
This lemma isolates the only non-local part of both arity-raising
constructions. It is formulated independently of their different
ways of orienting facets with three auxiliary vertices.

\begin{lemma}[Two-coordinate reflection]\label{lem:reflection}
Let $n\ge3$ and let $S$ be a set of $(n+1)$-vertex types.
Suppose $B_\ell\leftarrow A\rightarrow B_r$ is a one-point
diagram of $n$-hypertournaments in $\K^{(n)}_S$ without an ordinary
amalgam. Let $P$ be a set of auxiliary vertices, ordered before all
old vertices, such that
\[
 |P|>2^{\binom{|A|}{n-1}}.
\]
Suppose $C_\ell$ and $C_r$ are $(n+2)$-hypertournaments on
$P\cup B_\ell$ and $P\cup B_r$, agreeing on $P\cup A$,
whose signs satisfy
\begin{equation}\label{eq:two-pad-recovery}
 s_{C_i}(\{p,q\}\cup F)=s_{B_i}(F)
 \qquad(p<q\in P,\ F\in\tbinom{B_i}n,\ i\in\{\ell,r\}).
\end{equation}
Let
\[
 S^+=\{T_{n+2,d+1}:T_{n,d}\in S\}.
\]
If both $C_i$ belong to
$\K^{(n+2)}_{S^+\cup\{T_{n+2,0}\}}$, then they have no ordinary
amalgam in that class. The same conclusion holds, in particular,
when both belong to $\K^{(n+2)}_{S^+}$ and the target is that smaller class.
\end{lemma}
\begin{proof}
Suppose there is an ordinary amalgam. Restrict it to the union
of the two images. We first exclude identification of the new
vertices. Fix distinct $p,q\in P$ and choose their order as $p<q$.
The reconstruction must be interpreted \emph{intrinsically}:
on every injective ordered $n$-tuple $(v_1,\ldots,v_n)$ of old
vertices, take the orientation obtained by evaluating the
$(n+2)$-orientation on $(p,q,v_1,\ldots,v_n)$.
This defines an alternating $n$-orientation because permuting the
$v_i$ changes the larger orientation by the same parity. Equation
\eqref{eq:two-pad-recovery} guarantees that the reconstructed
structure restricts to both original sides. This argument is valid
even if $\ell$ and $r$ are identified: in that case no auxiliary
ordering of the amalgam need extend both side orderings
$\ell<A$ and $A<r$, but the intrinsic reconstruction does not
require one. It would give an ordinary amalgam of the source
diagram, a contradiction. Thus the new vertices are distinct.

We may now choose one auxiliary order on the distinct-vertex union
with $P<\ell<A<r$ and perform the remaining packet computation.

For each $p\in P$, record its \emph{one-pad mixed profile}
\[
 \chi(p)=\bigl(s_C(\{p,\ell,r\}\cup U)\bigr)_
 {U\in\binom A{n-1}}\in\{\pm1\}^{\binom{|A|}{n-1}}.
\]
By pigeonhole, there are $p<q$ with $\chi(p)=\chi(q)$.
Define an $n$-hypertournament $D$ on $A\cup\{\ell,r\}$ by
$s_D(F)=s_C(\{p,q\}\cup F)$. The recovery condition implies
that $D$ extends both old sides. For each mixed old
$(n+1)$-set $Y=\{\ell,r\}\cup U$ with
$U\in\binom A{n-1}$, form $X=\{p,q\}\cup Y$.
Its first two normalized facet signs are opposite because the
one-pad profiles agree; its other normalized facet signs are
exactly those of $D[Y]$, with indices shifted by two. Hence
\[
 d_{n+2}(X)=1+d_n(D[Y]).
\]
As $X$ is admissible and its type cannot be $T_{n+2,0}$,
we have $d_n(D[Y])\in S$. Every other $(n+1)$-set of $D$
lies in an old side, so $D\in\K^{(n)}_S$, a contradiction.
The same argument applies to the smaller target class.
\end{proof}

\section{The two-coordinate suspension in even arity}\label{sec:even-lift}
\begin{lemma}[Local type map]\label{lem:even-types}
Let $n=2k\ge4$ and let $B$ be an $n$-hypertournament.  Take a disjoint
linearly ordered set $P$, put all vertices of $P$ before $B$, and define an
$(n+2)$-hypertournament $B^{\uparrow P}$ on $P\cup B$ as follows:
\begin{itemize}
\item If an $(n+2)$-set $F$ contains exactly two vertices of $P$, give
  it the sign of the $n$-set $F\setminus P$ in $B$.
\item Otherwise give $F$ positive sign (increasing orientation).
\end{itemize}
Every $(n+3)$-set $X$ in $B^{\uparrow P}$ has central type
$T_{n+2,k+1}$, except possibly if $|X\cap P|=2$.  In that case
\[
  d_{n+2}(X)=d_n(X\setminus P)+1.
\]
Consequently, if $B\in\K^{(n)}_S$ and $T_{n,k}\in S$, then
\[
 B^{\uparrow P}\in\K^{(n+2)}_{S^+},\qquad
 S^+=\{T_{n+2,d+1}:T_{n,d}\in S\}.
\]
\end{lemma}
\begin{proof}
Write $h=|X\cap P|$.  If $h=0,1$ or $h\ge4$, every $(n+2)$-facet
of $X$ has positive sign, and the alternating normalized vector on the
odd number $n+3=2k+3$ of facets has $k+1$ or $k+2$ positive coordinates.
Hence $X$ is central.

If $h=3$, write the three auxiliary vertices first.  The normalized
coordinates arising by deletion of an auxiliary vertex are
$(s,-s,s)$, where $s$ is the sign of the common $n$-set of old vertices.
Exactly one or two of these coordinates are positive.  The remaining
$n=2k$ facets have positive raw signs and positions $3,\ldots,n+2$;
exactly $k$ of their normalized signs are positive.  Again the total
number of positive normalized signs is $k+1$ or $k+2$.

Finally, if $h=2$, the two auxiliary vertices are first.  Their
deletions produce normalized signs $(+,-)$, while deletion of any old
vertex produces precisely the corresponding normalized sign of the
old $(n+1)$-set (the index is shifted by two).  Hence the new positive
count is exactly one more than the old positive count.  Since the total
numbers of facets are $n+3$ and $n+1$, respectively, the defect increases
by one.
\end{proof}

\begin{proposition}[Lifting one-point non-amalgamation]\label{prop:even-lift}
Let $n=2k\ge4$ and let $T_{n,k}\in S$.
Suppose that a one-point diagram
\[
 B_\ell=A\cup\{\ell\}\longleftarrow A
 \longrightarrow B_r=A\cup\{r\}
\]
in $\K^{(n)}_S$ has no \emph{ordinary} amalgam.  Then the same is true
in arity $n+2$ for each of
\[
 \K^{(n+2)}_{S^+}
 \qquad\text{and}\qquad
 \K^{(n+2)}_{S^+\cup\{T_{n+2,0}\}}.
\]
In both classes the obstruction is again a one-point diagram.
\end{proposition}
\begin{proof}
Take the same ordered auxiliary set $P$ on both sides, with
$|P|>2^{\binom{|A|}{n-1}}$, and apply
Lemma~\ref{lem:even-types}.
The two constructions agree on the lifted common base $P\cup A$
and belong to $\K^{(n+2)}_{S^+}$; they satisfy the exact
recovery rule~\eqref{eq:two-pad-recovery} by definition.
Apply Lemma~\ref{lem:reflection}, first with target $S^+$ and
then with target $S^+\cup\{T_{n+2,0}\}$.
\end{proof}

\begin{remark}[A smaller auxiliary set from an unsatisfiable core]
\label{rem:core}
The profile need only record those $(n-1)$-sets $U$ for which the
corresponding mixed $(n+1)$-set is used in an unsatisfiability certificate
for the original one-point diagram.  If an unsatisfiable subsystem uses
$h$ such sets, then $|P|=2^h+1$ already suffices.
\end{remark}



\section{The rooted two-coordinate suspension in odd arity}\label{sec:odd-lift}
For odd $n=2q-1\ge3$, let $S$ contain $T_{n,q}$ and define
\[
 S^+=\{T_{n+2,d+1}:T_{n,d}\in S\}.
\]
Let $B\in\K^{(n)}_S$ contain a distinguished vertex $a$.
Choose a linear order on $B$ with $a$ first. Define on its $(n-1)$-sets
\[
 y_B(F)=
 \begin{cases}
 +1,&a\in F,\\
 s_B(\{a\}\cup F),&a\notin F.
 \end{cases}
 \tag{2.1}\label{eq:root-potential}
\]
This is well-defined because $|\{a\}\cup F|=n$ in the second case.

\begin{lemma}[Root identity]\label{lem:odd-root}
For each $F\in\binom Bn$,
\[
 s_B(F)-\partial y_B(F)=
 \begin{cases}
 0,&a\in F,\\
 \partial s_B(\{a\}\cup F),&a\notin F.
 \end{cases}
 \tag{2.2}\label{eq:root-identity}
\]
\end{lemma}
\begin{proof}
If $a\notin F$, expand $\partial s_B(\{a\}\cup F)$ along its
first vertex $a$; the remaining terms are exactly $-\partial y_B(F)$.
If $F=\{a\}\cup G$ then deleting $a$ contributes $s_B(F)$ to
$\partial y_B(F)$, while the remaining $n-1$ facets have $y_B=+1$.
Since $n-1$ is even, their alternating signed sum is zero.
\end{proof}

Fix an ordered auxiliary set $P$ disjoint from $B$ with $P<a$.
Define the sign of every $(n+2)$-set $Z\subseteq P\cup B$ by
\[
 s_{B^{\uparrow P}}(Z)=
 \begin{cases}
 s_B(Z\setminus P),& |Z\cap P|=2,\\
 y_B(Z\setminus P),& |Z\cap P|=3,\\
 +1,& |Z\cap P|\notin\{2,3\}.
 \end{cases}
 \tag{2.3}\label{eq:padding}
\]
Here $Z\setminus P$ has size $n$ or $n-1$ in the first two cases.

\begin{lemma}[Local type control]\label{lem:odd-types}
Let $X$ be an $(n+3)$-set in $B^{\uparrow P}$, and put
$h=|X\cap P|$. Every local type of $B^{\uparrow P}$ belongs to $S^+$.
More precisely,
\begin{enumerate}
\item If $h\in\{0,1\}$ or $h\ge4$, the type is central $T_{n+2,q+1}$.
\item If $h=2$, the type is $T_{n+2,d+1}$, where the old
$(n+1)$-set $X\setminus P$ has type $T_{n,d}$.
\item If $h=3$ and $a\notin X\setminus P$, the type is $T_{n+2,d+1}$,
where $\{a\}\cup(X\setminus P)$ has old type $T_{n,d}$.
If $h=3$ and $a\in X\setminus P$, the new type is central.
\end{enumerate}
\end{lemma}
\begin{proof}
The new normalized sign vector has length $n+3=2q+2$, and its
type is determined by the absolute value of the alternating sum of its
raw facet signs. We check the number $h$.

For $h=0,1$ or $h\ge5$, all facet signs are $+1$, so the alternating
sum is zero. If $h=2$, the two auxiliary deletions have raw sign $+1$
and cancel in the alternating sum. The remaining $n+1$ terms are the
old facet signs of $Y=X\setminus P$, with the same alternating
coefficients, so the sum is $\partial s_B(Y)$.

If $h=3$, the three auxiliary deletions have identical raw sign
$s_B(F)$, for $F=X\setminus P$, and their alternating sum is $s_B(F)$.
The remaining old deletions, whose positions begin at $3$, contribute
$-\partial y_B(F)$. Thus the total is
$s_B(F)-\partial y_B(F)$, described by Lemma~\ref{lem:odd-root}.

If $h=4$, the four auxiliary deletions contribute four alternating
copies of $y_B(F)$ and cancel. The remaining $n-1$ old deletions
have raw sign $+1$ and contribute zero because $n-1$ is even.

Consequently the absolute alternating sum is either zero or the
absolute alternating sum of an old $(n+1)$-set. Since the number of
normalized facets has increased by two, the local type index
increases by exactly one in the latter case. Central $T_{n+2,q+1}$
corresponds to old central $T_{n,q}$, which belongs to $S$.
\end{proof}

\begin{theorem}[Odd-arity obstruction transfer]\label{thm:odd-lift}
Let $n\ge3$ be odd, let $T_{n,q}\in S$, and suppose
$\K^{(n)}_S$ contains a one-point-versus-one-point diagram
\[
 B_\ell=A\cup\{\ell\}\longleftarrow A
 \longrightarrow B_r=A\cup\{r\}
\]
without any ordinary amalgam. Then both
\[
 \K^{(n+2)}_{S^+}
 \quad\text{and}\quad
 \K^{(n+2)}_{S^+\cup\{T_{n+2,0}\}}
\]
lack ordinary AP, witnessed by one-point diagrams.
\end{theorem}
\begin{proof}
The base $A$ can be taken nonempty (a one-point diagram with empty
base already has an amalgam in any unbounded hypertournament class
of arity at least three). Choose $a\in A$ and one order on $A$
with $a$ first; extend it to compatible orders on both old sides.
The rooted functions~\eqref{eq:root-potential} agree on every
$(n-1)$-subset of $A$, since the sides agree on $A$.
Choose an auxiliary $P$ with $|P|>2^{\binom{|A|}{n-1}}$.
The constructions~\eqref{eq:padding} therefore agree on $P\cup A$,
are members of $\K^{(n+2)}_{S^+}$ by
Lemma~\ref{lem:odd-types}, and satisfy the two-pad recovery
rule~\eqref{eq:two-pad-recovery}. Apply
Lemma~\ref{lem:reflection} to either target class.
\end{proof}

\begin{remark}[Unified arity raising]
The preceding theorem covers \emph{odd} arities. The unrooted padding of Lemma~\ref{lem:even-types}
covers even arities, while the rooted padding here covers odd ones.
Both transfer statements are instances of Lemma~\ref{lem:reflection}. The two constructions
are genuinely different: the root potential in \eqref{eq:root-potential}
compensates for the middle (three-pad) case that is problematic for odd $n$.
\end{remark}



\section{The homogeneous-free switching class in arities $3\pmod4$}\label{sec:three-four}
An $n$-hypertournament assigns an orientation to each $n$-subset.
For an auxiliary order, encode orientations by bits
$e:\binom Vn\to\F$. On an ordered $(n+1)$-set $X$, its facet packet is
$\bigl(e(X\setminus\{x_i\})\bigr)_{i=0}^{n}$.
The homogeneous type $T_{n,0}=H_{n+1}$ is exactly either of the two
alternating packets $(i\bmod2)_i$ or $(1+i\bmod2)_i$.
For $n\equiv3\pmod4$, the switching class is precisely
\[
\E_n=\left\{e:\sum_{x\in X}e(X\setminus\{x\})=0
\quad(X\in\tbinom V{n+1})\right\},
\]
whose permitted isomorphism types are $T_{n,0},T_{n,2},\ldots,
T_{n,(n+1)/2}$. The latter follows by comparing the parity of raw
facet bits with the parity of normalized orientations.
Write $\delta$ for the unsigned simplicial coboundary over $\F$:
\[
 (\delta z)(Y)=\sum_{Z\in\binom Y{|Y|-1}}z(Z).
\]
The equality $\E_n=\{\delta z:z:\binom V{n-1}\to\F\}$
follows from Lemma~\ref{lem:exactness}; it uses no external
cohomological input.

\begin{theorem}\label{thm:p3}
For every $n\ge7$ with $n\equiv3\pmod4$,
\[
 \boxed{\K^{(n)}_{\{T_{n,2},T_{n,4},\ldots,T_{n,(n+1)/2}\}}
 \text{ does not have the amalgamation property.}}
\]
In particular, $\K^{(7)}_{\{T_{7,2},T_{7,4}\}}$ does not
have ordinary amalgamation.
\end{theorem}

\subsection{An admissible random one-point diagram}
Put $m=n-1$, so $m\ge6$ and $m\equiv2\pmod4$. Choose $t\ge2$ and a
linearly ordered set
\[
 A=V_0< V_1<\cdots<V_{m-1},\qquad |V_i|=t.
\]
Adjoin $\ell<A<r$. Split the block indices into consecutive halves
\[
 I=\{0,\ldots,m/2-1\},\qquad
 J=\{m/2,\ldots,m-1\}.
\]
We prescribe an $(n-1)=m$-ary potential $z_A$ on $A$ and two
extensions of it, $z_\ell$ on $A\cup\{\ell\}$ and
$z_r$ on $A\cup\{r\}$:
\begin{itemize}
\item On $m$-sets of $A$ meeting every block exactly once, set
$z_A(Q)=W(Q)$, independently and uniformly in $\F$; on all other
$m$-sets put $z_A=0$.
\item On $m$-sets containing $\ell$, set
$z_\ell(\{\ell\}\cup U)=L_i(U)$ if $U$ is transversal to precisely
those blocks other than $V_i$, with $i\in I$. All variables $L_i(U)$
are independent uniform bits; set all remaining such potential values
to zero. The restriction of $z_\ell$ to $A$ equals $z_A$.
\item On $m$-sets containing $r$, analogously set
$z_r(U\cup\{r\})=R_j(U)$ when $U$ is transversal to exactly the
blocks other than $V_j$, for $j\in J$. These are fresh independent
uniform bits; put zero elsewhere, and keep $z_A$ on $A$.
\end{itemize}
All families $W,L,R$ are mutually independent. Orient the $n$-sets of
the two sides using
$e_\ell=\delta z_\ell$ and $e_r=\delta z_r$.
The structures agree on the common induced substructure $A$.

\begin{lemma}[Protected zero facets]\label{lem:p3-admissible}
For \emph{every} assignment of all the random potential bits, neither
$e_\ell$ nor $e_r$ has a homogeneous $(n+1)$-set.
Thus these are members of the class in Theorem~\ref{thm:p3}.
\end{lemma}
\begin{proof}
Each side belongs to $\E_n$, so it suffices to exclude the two
alternating facet packets. We show that every $(m+2)$-set in either
side has two facets with bit zero at indices of opposite parity.

First take $X\subseteq A$ with $|X|=m+2$. If it occupies fewer than
$m$ blocks, every facet bit is zero: no potential value supported on
$A$ is available. If $X$ meets all $m$ blocks, there are only two
extra points beyond a transversal. Hence at most two blocks are
non-singleton. Every non-singleton block contributes either an
adjacent pair, or one block contributes an adjacent triple. In either
case these points occupy at most two positions of each parity in
$X$. There are $m/2+1\ge4$ positions of each parity, so at least
one singleton-block vertex occurs in each parity. Deleting such a
vertex removes its entire block and forces its facet bit to be zero.

Next take $X=\{\ell\}\cup Y$, $|Y|=m+1$. If $Y$ meets at most
$m-2$ blocks, all relevant potentials are zero. If $Y$ meets exactly
$m-1$ blocks, there are two extra points, distributed as either two
adjacent pairs or one adjacent triple. As $m\ge6$, among the
singleton-block vertices of $Y$ are positions of both parities.
Deleting any such vertex removes a block, leaving at most $m-2$
distinct base blocks; neither a base potential supported on all $m$
blocks nor a left potential supported on $m-1$ blocks can survive.

Finally suppose $Y$ meets all $m$ blocks. Exactly one block is
represented twice, and its vertices are adjacent in the order.
Among the last $m/2\ge3$ blocks (the set $J$), at least two
singleton blocks occur in positions of opposite parity: if the
doubled block lies outside $J$, all $J$ positions are consecutive;
if it lies inside $J$, deleting the adjacent doubled pair still leaves
singleton positions of both parities. Deleting either protected
singleton loses a block of $J$, so every permitted left potential
(which uses all blocks in $J$) and every base potential vanish on
that facet.

The case $X=Y\cup\{r\}$ is symmetric, using singleton blocks from
$I$. Thus every possible packet has two fixed zero coordinates of
opposite parity, excluding both alternating patterns.
\end{proof}

\subsection{All possible mixed completions}
Put $h(U)=z_\ell(\{\ell\}\cup U)+z_r(U\cup\{r\})$ for
$U\in\binom A{m-1}$. Let
$\alpha(Q)=e_\ell(\{\ell\}\cup Q)$ and
$\beta(Q)=e_r(Q\cup\{r\})$ for $Q\in\binom Am$.
Directly from the potential formula,
\begin{equation}\label{eq:links}
 \alpha(Q)+\beta(Q)=(\delta h)(Q).
\end{equation}
Suppose an $n$-hypertournament on
$A\cup\{\ell,r\}$ extends both sides and satisfies the switching
parity constraints. Encode its mixed $n$-sets by
\[
 x(U)=e(\{\ell,r\}\cup U),\qquad U\in\tbinom A{m-1}.
\]
On each mixed $(n+1)$-set the parity constraint gives
$\delta x=\alpha+\beta=\delta h$.
By Lemma~\ref{lem:exactness}, applied to the cocycle $x+h$,
there is a cochain $u:\binom A{m-2}\to\F$ such that
\begin{equation}\label{eq:gauge}
 \boxed{x=h+\delta u.}
\end{equation}
This covers \emph{all} possible mixed completions satisfying parity,
not just completions arising from a preferred potential.

A \emph{transversal} $Q=(v_0,\ldots,v_{m-1})$, $v_i\in V_i$, has
$m$ mixed facet bits
\[
 x_i(Q)=x(Q\setminus\{v_i\}),\qquad 0\le i<m.
\]
For fixed $u$, these are independent unbiased bits, since
\[
 h(Q\setminus\{v_i\})=
 \begin{cases}
 L_i(Q\setminus\{v_i\}),&i\in I,\\
 R_i(Q\setminus\{v_i\}),&i\in J,
 \end{cases}
\]
and $\delta u$ is fixed.
Call a transversal $Q$ \emph{alternating for $u$} if
$(x_i(Q))_{i<m}$ is either alternating vector of length $m$.
Let $M_u$ count such transversals. There are $t^m$ transversals, and
\begin{equation}\label{eq:mean}
 \mathbb E[M_u]=p t^m,\qquad p=\frac{2}{2^m}=2^{1-m}.
\end{equation}

\begin{lemma}[Independent mixed polarity]\label{lem:p3-polarity}
Condition on all the variables $L,R$ and on a fixed gauge $u$.
For each transversal $Q$ alternating for $u$, exactly one of the two
values of the still-independent bit $W(Q)$ makes
$\{\ell,r\}\cup Q$ a forbidden homogeneous $(n+1)$-set.
Consequently the conditional probability that all transversal mixed
$(n+1)$-sets are admissible is at most $2^{-M_u}$.
\end{lemma}
\begin{proof}
The facet packet on $\{\ell,r\}\cup Q$ is
\[
 \bigl(\beta(Q),x_0(Q),\ldots,x_{m-1}(Q),\alpha(Q)\bigr).
\]
As $m\equiv2\pmod4$, any alternating interior $m$-vector has odd
parity. Thus $\alpha(Q)+\beta(Q)=\sum_i x_i(Q)=1$.
If the interior is alternating and the two endpoints are opposite,
there are exactly two possible endpoint pairs, one of which produces
an alternating packet of length $m+2$ and one which does not.
But
\[
 \alpha(Q)=W(Q)+\sum_{i\in I}L_i(Q\setminus\{v_i\}),
 \qquad
 \beta(Q)=W(Q)+\sum_{j\in J}R_j(Q\setminus\{v_j\}).
\]
Toggling $W(Q)$ toggles both endpoints, so one choice is forbidden.
The $W(Q)$ for different transversals are independent, which proves
the bound.
\end{proof}

\subsection{Independent transversals and the union bound}
Identify each block $V_i$ with the cyclic group $\mathbb Z/t\mathbb Z$.
Let $\mathcal Q$ consist of those transversals
$Q=(v_0,\ldots,v_{m-1})$ whose block coordinates sum to zero:
\[
 \mathcal Q=\{Q: v_i\in\mathbb Z/t\mathbb Z,\quad
                  v_0+\cdots+v_{m-1}=0\}.
\]
There are $t^{m-1}$ such transversals. No two distinct members
of $\mathcal Q$ agree in $m-1$ coordinates, since the missing coordinate
is determined by the sum condition. Consequently their
$(m-1)$-facets are all distinct.

Fix a gauge $u$. For every $Q\in\mathcal Q$, the $m$ values
$h(Q\setminus\{v_i\})$ are independent uniform bits. Moreover, by
the distance-two property, the sets of random source bits used by
different members of $\mathcal Q$ are disjoint. Since $\delta u$
is fixed, the $m$ interior bits $x_i(Q)$ are independent and uniform,
independently for all $Q\in\mathcal Q$.
Thus the probability that the interior $m$-vector is alternating is
$2/2^m=2^{1-m}$. Conditional on that event, Lemma~\ref{lem:p3-polarity}
says that the independent base bit $W(Q)$ makes the whole mixed packet
forbidden with probability $1/2$. The events of getting a forbidden
homogeneous $(n+1)$-set on $\{\ell,r\}\cup Q$ are therefore
\emph{independent}, each with probability $2^{-m}$. In particular,
\begin{equation}\label{eq:fixed}
 \Pr[\text{all selected mixed packets are admissible for a fixed }u]
 = (1-2^{-m})^{t^{m-1}}.
\end{equation}
There are $2^{\binom{mt}{m-2}}$ possible cochains $u$. Any strong
amalgam in the desired class would give one of these gauges and would
avoid every selected forbidden packet. Therefore the union bound gives
\begin{equation}\label{eq:total}
 \Pr[\text{some parity-respecting strong amalgam exists}]
 \le 2^{\binom{mt}{m-2}}(1-2^{-m})^{t^{m-1}}.
\end{equation}
We can choose $t$ explicitly, avoiding any asymptotic or
analytic estimate. Put $M=2^m$ and
\[
 t=2^{m+1}m^{m-2},\qquad
 K=\frac{t^{m-1}}M=2(mt)^{m-2}.
\]
The elementary Bernoulli inequality yields
\[
 (1-1/M)^M=(1+1/(M-1))^{-M}
 \le (1+M/(M-1))^{-1}<\tfrac12.
\]
Partition the $t^{m-1}=MK$ selected transversals into $K$
blocks of size $M$. Since
$\binom{mt}{m-2}\le(mt)^{m-2}$, the right-hand side of
\eqref{eq:total} is at most
\[
 2^{\binom{mt}{m-2}-K}
 \le2^{-(mt)^{m-2}}<\tfrac12.
\]
This is a completely finite inequality suitable for a counting
formalisation. (The displayed $t$ is a generous bound, not an
attempt to minimise the witness size.)

To also exclude identification of $\ell$ and $r$, fix one transversal
$Q_0$. The event $\alpha(Q_0)\ne\beta(Q_0)$ has probability
$1/2$, independently of the base potential $W$. Since the
probability on the left of \eqref{eq:total} is less than $1/2$,
some assignments simultaneously satisfy $\alpha(Q_0)\ne\beta(Q_0)$
and admit no strong amalgam. If $\ell,r$ were identified in an
ordinary amalgam, their orientations on the $n$-sets extending
$Q_0$ would coincide; there is no sign discrepancy on moving the
extra vertex across $m=n-1$ base vertices because $m$ is even.
This contradicts $\alpha(Q_0)\ne\beta(Q_0)$ and completes the
proof of Theorem~\ref{thm:p3}.



\section{The critical classes in arities $1\pmod4$}\label{sec:one-four}

Fix an auxiliary linear order. We encode an oriented $n$-set $F$ by
$e(F)\in\F$, where $e(F)=0$ means its increasing enumeration represents
the chosen orientation. For an ordered $(n+1)$-set
$X=\{x_0<\cdots<x_n\}$, its \emph{raw facet packet} is
\[
 \bigl(e(X\setminus\{x_i\})\bigr)_{i=0}^{n}.
\]
The homogeneous type $T_{n,0}=H_{n+1}$ has precisely the two packets
$(i\bmod2)_{i=0}^{n}$ and $(1+i\bmod2)_{i=0}^{n}$. Write
$T_{n,d}$ for the type with $d$ minority normalized signs. Let
$N=n+1=4k+6$, so $N/2=2k+3$ is odd. The parity of the raw packet equals
$d+N/2\pmod2$; in particular, the condition
\[
  T_{n,d}(X)\in\{T_{n,0},T_{n,1},T_{n,3},\ldots,T_{n,N/2}\}
\]
is equivalent to
\begin{equation}\label{eq:admissible}
 \sum_{F\in\binom{X}{n}}e(F)=0
 \quad\text{or}\quad
 \bigl(e(X\setminus\{x_i\})\bigr)_{i=0}^{N-1}
 \text{ is alternating}.
\end{equation}
Denote this class by $\mathcal P_n$. Its homogeneous-set indicator is
\[
 D(X)=\sum_{F\in\binom Xn}e(F)\in\F,
 \qquad X\in\binom V N.
\]
By~\eqref{eq:admissible}, $D(X)=1$ exactly when $X$ is homogeneous.

\subsection{A uniform one-point diagram}

Let $m=4k$, and take the ordered common base
\[
 A=P\cup U,\qquad
 P=\{p_0<\cdots<p_{m-1}\}<
 U=\{a_0<\cdots<a_5\}.
\]
Thus $|A|=m+6=N$. For $k=0$, $P$ is empty. Put the new vertices in
one auxiliary order as
\[
 P<\ell<U<r.
\]
We prescribe two $n$-hypertournaments on $A\cup\{\ell\}$ and
$A\cup\{r\}$, agreeing on $A$. Since $n=N-1$, every $n$-set in
$A$ is $A\setminus\{v\}$ for some $v\in A$; every $n$-set
containing either new vertex is indexed by a pair of omitted base
vertices.

\paragraph{The common base.}
Set
\begin{equation}\label{eq:base}
 b_v:=e(A\setminus\{v\})=
 \begin{cases}
   1,&v\in P,\\
   (1,1,0,1,0,1)_j,&v=a_j.
 \end{cases}
\end{equation}
The raw packet on $A$ has parity $m+4=0\pmod2$ and is admissible.

\paragraph{The two link graphs.}
For $v,w\in A$, $v<w$, set
\[
 \alpha_{vw}=e(\{\ell\}\cup(A\setminus\{v,w\})),
 \qquad
 \beta_{vw}=e(\{r\}\cup(A\setminus\{v,w\})).
\]
The following data give every link bit:
\begin{enumerate}
\item For $v=a_i,w=a_j$ in $U$ ($i<j$), let the zero sets be
\begin{align}\label{eq:oldedges}
  \alpha_{a_i a_j}=0&\iff ij\in\{14,15,24,34,35\},\\
  \beta_{a_i a_j}=0&\iff ij\in\{01,03,05,14,25,34\}.
\end{align}
All unlisted core--core bits are $1$. Here $ij$ abbreviates the
unordered pair $\{i,j\}$.
\item For $p_s\in P$ and $a_j\in U$, put
\begin{equation}\label{eq:paddededges}
 \beta_{p_s a_j}=
 \begin{cases}
  s\bmod2,&j=0,\\
  1+s\bmod2,&j=4,\\
  1,&j\in\{1,2,3,5\},
 \end{cases}
 \qquad
 \alpha_{p_s a_j}=1+\beta_{p_s a_j}\pmod2.
\end{equation}
\item For two distinct pads $p_s,p_t\in P$, put
$\alpha_{p_sp_t}=\beta_{p_sp_t}=0$.
\end{enumerate}

\begin{lemma}[Admissibility of both sides]\label{lem:p1-sides}
The preceding prescriptions define members of $\mathcal P_n$.
Among their $N$-vertex substructures, the left side has no
homogeneous one; the right side has precisely the two homogeneous
sets
\[
 (A\setminus\{a_0\})\cup\{r\},\qquad
 (A\setminus\{a_4\})\cup\{r\}.
\]
\end{lemma}
\begin{proof}
The base $A$ is already admissible by~\eqref{eq:base}.
Fix a base vertex $v$. On the $N$-set
$(A\setminus\{v\})\cup\{\ell\}$, the raw packet consists of
$b_v$ and the link bits $\alpha_{vw}$ with $w\in A\setminus\{v\}$;
the corresponding right-side packet uses $\beta$. Hence their
parities are
\[
 b_v+\sum_{w\ne v}\alpha_{vw},\qquad
 b_v+\sum_{w\ne v}\beta_{vw}.
\]
For the six core vertices, the data~\eqref{eq:base} and
\eqref{eq:oldedges} give
\[
 \sum_{\substack{w\in U\\w\ne a_j}}\alpha_{a_jw}=b_{a_j},
 \qquad
 \sum_{\substack{w\in U\\w\ne a_j}}\beta_{a_jw}
 =b_{a_j}+\mathbf1_{j\in\{0,4\}}
 \quad(\bmod2).
\]
All pad contributions to these sums vanish: for each core vertex,
its prescribed pad row contains either $m$ identical bits or $m/2$
ones, both even since $m=4k$.

For a pad vertex $p_s$, the other-pad entries are zero; the six
pad--core link bits have parity $1$ on each side, because the $j=0$
and $j=4$ bits have opposite values, and the other four contribute
an even number. This equals $b_{p_s}$.

Therefore all left packets have even parity, and all right packets
except the two indexed by $a_0,a_4$ have even parity. For those
remaining two right packets, the full raw vectors, including their
$m$ pad-deletion coordinates, are respectively
\[
 \underbrace{0101\cdots01}_{m\text{ bits}},\ 010101,
 \qquad
 \underbrace{1010\cdots10}_{m\text{ bits}},\ 101010
\]
(where for $m=0$ the initial strings are empty). Both are
alternating, as required by~\eqref{eq:admissible}.
\end{proof}

\subsection{The forced contradiction}

\begin{theorem}\label{thm:p1}
For every $n\geq5$ satisfying $n\equiv1\pmod4$, the class
\[
 \boxed{\mathcal P_n
 =\K^{(n)}_{\{T_{n,0},T_{n,1},T_{n,3},\ldots,T_{n,(n+1)/2}\}}}
\]
does not have the ordinary amalgamation property.
\end{theorem}
\begin{proof}
Use the one-point diagram from Lemma~\ref{lem:p1-sides}.
Suppose first that it had a strong amalgam on
$V=A\cup\{\ell,r\}$, with $|V|=N+2$.

Form a graph $G$ on $V$ by declaring $vw$ an edge if and only if
$V\setminus\{v,w\}$ is a homogeneous $N$-set. By
Lemma~\ref{lem:zero-two}, every vertex of $G$ has degree either $0$ or $2$.
Lemma~\ref{lem:p1-sides} gives the fixed edges
\begin{equation}\label{eq:fixed-edges}
 \ell a_0,\qquad \ell a_4,
\end{equation}
no edges incident with $r$, and no edge $\ell r$.

Consider a potential edge $vw$ with $v,w\in A$. The associated
mixed $N$-set $X=V\setminus\{v,w\}$ contains both $\ell$ and $r$.
In the increasing facet packet of $X$, deleting $\ell$ gives the
right link bit $\beta_{vw}$, while deleting $r$ gives the left link
bit $\alpha_{vw}$. The position of $r$ is $N-1=m+5$, which is odd;
the position of $\ell$ is the number of pads retained in $X$.
Consequently, alternating packets are possible only in the cases
\begin{center}
\begin{tabular}{@{}lll@{}}
\toprule
Types of $v,w$ & Position of $\ell$ & Necessary endpoint relation\\
\midrule
both pads & $m-2$ (even) & $\alpha_{vw}\ne\beta_{vw}$\\
one pad, one core vertex & $m-1$ (odd) & $\alpha_{vw}=\beta_{vw}$\\
both core vertices & $m$ (even) & $\alpha_{vw}\ne\beta_{vw}$\\
\bottomrule
\end{tabular}
\end{center}
By construction, neither of the first two cases can be homogeneous:
both-pad links agree, and pad--core links disagree. For two core
vertices, $\alpha_{a_i a_j}\ne\beta_{a_i a_j}$ exactly on
\begin{equation}\label{eq:eligible}
 01,\ 03,\ 05,\ 15,\ 24,\ 25,\ 35.
\end{equation}
Thus these seven core pairs are the \emph{only} possible mixed edges of $G$.

The fixed edge $\ell a_4$ means $a_4$ already has degree at least one.
By the zero-or-two lemma its degree must be exactly two, and by
\eqref{eq:eligible} its only possible second neighbor is $a_2$.
So $a_2a_4\in G$. The vertex $a_2$ has no fixed edge to a new
vertex, and its only two eligible mixed neighbors are $a_4,a_5$.
Its degree must likewise be two, forcing $a_2a_5\in G$.

Now examine the common $n$-set
\[
 F=\{\ell,r\}\cup\bigl(A\setminus\{a_2,a_4,a_5\}\bigr).
\]
The homogeneous set $X_{24}=V\setminus\{a_2,a_4\}$ has right
endpoint bit $\beta_{a_2a_4}=1$ at the even packet position $m$.
Its facet obtained by deleting $a_5$ is exactly $F$, at packet
position $m+4$ (also even), so homogeneity forces
\[
 e(F)=1.
\]
On the other hand, the homogeneous set
$X_{25}=V\setminus\{a_2,a_5\}$ has right endpoint bit
$\beta_{a_2a_5}=0$ at position $m$, and its facet obtained by
deleting $a_4$ is again $F$ at position $m+4$. Thus homogeneity
forces $e(F)=0$, a contradiction.

Finally, the new vertices cannot be identified in an ordinary
amalgam: on the core pair $24$ the two link bits are
$\alpha_{a_2a_4}=0$ and $\beta_{a_2a_4}=1$. Identifying
$\ell$ with $r$ would move that vertex past the four retained
core vertices, an even permutation, so these bits would have to
agree. Hence there is no ordinary amalgam either.
\end{proof}

\begin{remark}[What has become of the old arity-five SAT seed?]
For $k=0$ the diagram is precisely the known six-vertex-base
arity-five obstruction for $\operatorname{Forb}(T_{5,2})$.
The present argument proves its non-amalgamability without SAT or
exhaustive enumeration. In its old coordinates, the contradictory
common mixed facet is the one labelled $x_{013}$.
The graph of homogeneous defects reveals the short reason why that
facet is forced to have both signs.
\end{remark}


\section{Deduction of the general-arity negative implications}\label{sec:conclusion}
\begin{proof}[Proof of Theorem~\ref{thm:classification}: negative direction]
We work only with unbounded classes, so the central type is present.
The separate positive constructions are not used as proof premises;
we now establish non-amalgamation for every unbounded set of types
excluded by Theorem~\ref{thm:classification}.

\emph{Even arities.}  The base arity $n=4$ is given by
Lemmas~\ref{lem:4-central} and~\ref{lem:4-forb}.
For $n=2k\ge4$, assume that every unbounded $n$-ary class missing
$T_{n,k-1}$ has a one-point diagram with no ordinary amalgam.
Consider $S'\subseteq\{T_{n+2,0},\ldots,T_{n+2,k+1}\}$
containing the central type but omitting $T_{n+2,k}$, and set
\[
 S=\{T_{n,d}:T_{n+2,d+1}\in S'\}.
\]
Then $S$ contains $T_{n,k}$ but omits $T_{n,k-1}$.
Proposition~\ref{prop:even-lift} transfers the negative diagram to
$S^+$ with or without the homogeneous type $T_{n+2,0}$,
which is exactly the class specified by $S'$.
This proves the theorem in every even arity.

\emph{Odd arities.}  At $n=5$ the three unbounded classes not
covered by a positive construction are
\[
 \{T_{5,3}\},\qquad \{T_{5,0},T_{5,3}\},\qquad
 \{T_{5,0},T_{5,1},T_{5,3}\}.
\]
They fail ordinary AP by Proposition~\ref{prop:cone} applied to the
arity-four central-only seed, by Lemma~\ref{lem:5-two}, and by
Theorem~\ref{thm:p1}, respectively.  Thus all three negative
arity-five cases are proved mathematically.

Now let $n\ge7$ be odd and put $q=(n+1)/2$.
Take an unbounded class $\K^{(n)}_{S'}$ which omits $T_{n,q-1}$,
and define $S=\{T_{n-2,d}:T_{n,d+1}\in S'\}$.
This is an unbounded set of types in arity $n-2$, missing
$T_{n-2,q-2}$.  If $S$ is not the switching set $E_{n-2}$,
the odd-arity induction hypothesis says that $\K^{(n-2)}_S$
fails ordinary AP.  The rooted transfer
Theorem~\ref{thm:odd-lift} then implies the same for the target
$S'$, independently of whether $T_{n,0}$ is included.

It remains to examine $S=E_{n-2}$.  Shifting its type indices up by
one gives all positive-index types of $E_n$; the parity calculation
shows the following:
\[
\begin{array}{c|c|c}
 n\bmod4&\text{switching set }E_n&\text{other set with the same positive indices}\\\hline
 3&\{T_{n,0},T_{n,2},\ldots,T_{n,q}\}&
 E_n\setminus\{T_{n,0}\}\\
 1&\{T_{n,1},T_{n,3},\ldots,T_{n,q}\}&
 E_n\cup\{T_{n,0}\}.
\end{array}
\]
The switching case is outside the scope of this negative implication.
The other case fails ordinary AP by Theorem~\ref{thm:p3} when
$n\equiv3\pmod4$ and by Theorem~\ref{thm:p1} when
$n\equiv1\pmod4$.
This closes the induction for every odd arity at least five.
\end{proof}

\begin{remark}[Scope and proof status]
The general-arity argument in Sections~\ref{sec:arity4}--
\ref{sec:conclusion} is mathematical: it does not take a solver output,
an exhaustive electronic certificate, or a finite sample of larger
arities as a premise.  Two short finite tables below
make all seed calculations reproducible by direct substitutions and
implications.  The probabilistic proof is an existence argument and
may produce an enormous finite obstruction, but is not
computer-assisted.  That general-arity proof is adapted from the source manuscript and
has not yet received line-by-line independent mathematical review.
The separate five-constrained ternary classification is
computer-assisted and intentionally omitted from this note.
Nothing here has been formalised in Lean.
\end{remark}



\appendix

\clearpage
\section{Hand-checkable finite tables for the arity-four seed}
\label{app:four-tables}
For Table~\ref{tab:n4extend}, the six-bit word is written in the order
\[
 (v_{12},v_{13},v_{14},v_{23},v_{24},v_{34}),
\]
and the four-bit extensions in the order
$(v_{15},v_{25},v_{35},v_{45})$.  The six constraints involving
vertex $5$ are checked using the triple types $(*)$ above.
The elementary three-variable forbidden-pattern rule gives the
initial six assignments on $\{1,2,3\}$; extending to $4$
leaves precisely the sixteen rows listed here.

\begin{table}[htbp]\centering\small
\begin{tabular}{@{}cc|cc@{}}\toprule
Four-vertex word&Allowed four new bits&Four-vertex word&Allowed four new bits\\\midrule
\texttt{000000}&\texttt{1101}&\texttt{011110}&\texttt{--}\\
\texttt{100000}&\texttt{1101}&\texttt{111110}&\texttt{1011}\\
\texttt{101000}&\texttt{1100, 1101}&\texttt{110001}&\texttt{0010}\\
\texttt{111000}&\texttt{--}&\texttt{101001}&\texttt{1100, 0010}\\
\texttt{000100}&\texttt{0111}&\texttt{001011}&\texttt{1100}\\
\texttt{111100}&\texttt{0000, 1111}&\texttt{101011}&\texttt{1100}\\
\texttt{000110}&\texttt{0011}&\texttt{010111}&\texttt{0110, 0011}\\
\texttt{001110}&\texttt{--}&\texttt{001111}&\texttt{0110}\\
\bottomrule\end{tabular}
\caption{All admissible restrictions to vertices $1,2,3,4$ and all extensions to vertex $5$.  A dash means that no extension is allowed.}\label{tab:n4extend}
\end{table}

In Table~\ref{tab:n4paths}, the ten-bit word records, in order,
\[
 (v_{12},v_{13},v_{14},v_{15},v_{23},v_{24},v_{25},v_{34},v_{35},v_{45}).
\]
In an implication path, $i$ denotes $y_i=1$ and $\bar\imath$
denotes $y_i=0$.  Adjacent literals in a path have distinct indices;
the step is the elementary implication supplied by that pair in $(**)$.
Each row exhibits implications both ways between opposite literals,
so is unsatisfiable without any further assumptions.

{\small
\begin{longtable}{@{}cll@{}}
\caption{Explicit $2$-SAT contradictions for every possible restriction to $1,2,3,4,5$.\label{tab:n4paths}}\\
\toprule
Ten-bit word&Path $\neg y_i\Rightarrow y_i$&Path $y_i\Rightarrow\neg y_i$\\\midrule
\endfirsthead
\toprule Ten-bit word&Path $\neg y_i\Rightarrow y_i$&Path $y_i\Rightarrow\neg y_i$\\\midrule
\endhead
\texttt{1110100000}&$\bar{1}\Rightarrow2\Rightarrow3\Rightarrow4\Rightarrow5\Rightarrow1$&$1\Rightarrow\bar{2}\Rightarrow\bar{4}\Rightarrow\bar{1}$\\
\texttt{1011001000}&$\bar{1}\Rightarrow2\Rightarrow5\Rightarrow3\Rightarrow1$&$1\Rightarrow\bar{2}\Rightarrow\bar{4}\Rightarrow\bar{1}$\\
\texttt{1011001100}&$\bar{3}\Rightarrow4\Rightarrow5\Rightarrow3$&$3\Rightarrow1\Rightarrow\bar{2}\Rightarrow\bar{3}$\\
\texttt{0011011100}&$\bar{3}\Rightarrow4\Rightarrow5\Rightarrow3$&$3\Rightarrow1\Rightarrow4\Rightarrow\bar{3}$\\
\texttt{1011011100}&$\bar{3}\Rightarrow4\Rightarrow5\Rightarrow3$&$3\Rightarrow1\Rightarrow\bar{2}\Rightarrow\bar{3}$\\
\texttt{1100000110}&$\bar{5}\Rightarrow\bar{3}\Rightarrow4\Rightarrow5$&$5\Rightarrow1\Rightarrow\bar{2}\Rightarrow\bar{5}$\\
\texttt{1010000110}&$\bar{5}\Rightarrow\bar{3}\Rightarrow4\Rightarrow5$&$5\Rightarrow1\Rightarrow\bar{2}\Rightarrow\bar{5}$\\
\texttt{0100111110}&$\bar{1}\Rightarrow\bar{4}\Rightarrow3\Rightarrow5\Rightarrow1$&$1\Rightarrow2\Rightarrow4\Rightarrow\bar{3}\Rightarrow\bar{1}$\\
\texttt{0010111110}&$\bar{1}\Rightarrow\bar{3}\Rightarrow4\Rightarrow5\Rightarrow1$&$1\Rightarrow2\Rightarrow3\Rightarrow\bar{4}\Rightarrow\bar{1}$\\
\texttt{0001001001}&$\bar{4}\Rightarrow\bar{3}\Rightarrow\bar{5}\Rightarrow4$&$4\Rightarrow1\Rightarrow5\Rightarrow\bar{4}$\\
\texttt{1001001001}&$\bar{4}\Rightarrow\bar{3}\Rightarrow\bar{5}\Rightarrow4$&$4\Rightarrow1\Rightarrow\bar{2}\Rightarrow\bar{4}$\\
\texttt{1011001001}&$\bar{4}\Rightarrow\bar{3}\Rightarrow\bar{5}\Rightarrow4$&$4\Rightarrow2\Rightarrow5\Rightarrow\bar{4}$\\
\texttt{0000110011}&$\bar{2}\Rightarrow\bar{1}\Rightarrow\bar{4}\Rightarrow5\Rightarrow2$&$2\Rightarrow3\Rightarrow5\Rightarrow\bar{4}\Rightarrow\bar{2}$\\
\texttt{1111110011}&$\bar{3}\Rightarrow\bar{1}\Rightarrow2\Rightarrow3$&$3\Rightarrow4\Rightarrow\bar{5}\Rightarrow\bar{3}$\\
\texttt{0000101011}&$\bar{2}\Rightarrow\bar{1}\Rightarrow\bar{5}\Rightarrow4\Rightarrow2$&$2\Rightarrow3\Rightarrow4\Rightarrow\bar{5}\Rightarrow\bar{2}$\\
\texttt{1111101011}&$\bar{3}\Rightarrow\bar{1}\Rightarrow2\Rightarrow3$&$3\Rightarrow4\Rightarrow\bar{5}\Rightarrow\bar{3}$\\
\texttt{0100110111}&$\bar{4}\Rightarrow\bar{2}\Rightarrow\bar{5}\Rightarrow4$&$4\Rightarrow1\Rightarrow3\Rightarrow\bar{4}$\\
\bottomrule
\end{longtable}
}

\clearpage
\section{The sixteen homogeneous-defect supports in arity five}
\label{app:five-tables}
The $19$ possible missing triples are listed in~\eqref{eq:eligible5}.
Every solution of the pair-parity equations is uniquely described by
$S_*\triangle aU_a\triangle bU_b\triangle cU_c\triangle dU_d$
as in~\eqref{eq:four-free5}.  Table~\ref{tab:n5supports}
checks the \emph{integer} pair-degree upper bound in all sixteen
cases.  A violating pair has degree four (including the fixed
right-side contribution).  Thus only the five rows marked A--E survive.

\begin{table}[htbp]\centering\small
\begin{tabular}{@{}cccc|cccc@{}}\toprule
$(d,c,b,a)$&Status&&&$(d,c,b,a)$&Status&&\\\midrule
\texttt{0000}&OK (E)&&&\texttt{1000}&OK (C)&&\\
\texttt{0001}&$d(45)=4$&&&\texttt{1001}&OK (D)&&\\
\texttt{0010}&OK (B)&&&\texttt{1010}&$d(23)=4$&&\\
\texttt{0011}&$d(04)=4$&&&\texttt{1011}&$d(04)=4$&&\\
\texttt{0100}&$d(26)=4$&&&\texttt{1100}&OK (A)&&\\
\texttt{0101}&$d(14)=4$&&&\texttt{1101}&$d(14)=4$&&\\
\texttt{0110}&$d(36)=4$&&&\texttt{1110}&$d(23)=4$&&\\
\texttt{0111}&$d(04)=4$&&&\texttt{1111}&$d(04)=4$&&\\
\bottomrule
\end{tabular}
\caption{Complete degree test for the four free homogeneous-defect bits.}\label{tab:n5supports}
\end{table}


\section{Lean-oriented proof interfaces and dependency ledger}
\label{app:lean}
\paragraph{Definitions to formalise.}
An alternating orientation on injective ordered $n$-tuples,
induced structures, arbitrary vertex embeddings (\emph{not}
order-preserving embeddings), ordinary AP, normalized facet packets,
and the invariant $d_n$. The signed-permutation calculation in
Section~\ref{sec:statement} makes the relation between $d_n$ and
the intrinsic isomorphism type explicit.

The public theorem should be Corollary~\ref{cor:export}, rather than
the two parity-split negative implications: it is the exact necessary
half of the future iff classification and hides all seed-specific
choices. Allowed type sets must be well-formed (bounded by $q$).
A recommended formal signature is conceptually
\[
  n\ge4,\quad S\subseteq\{0,\ldots,q\},\quad
  q\in S,\quad\operatorname{AP}(\K^{(n)}_S)
  \Longrightarrow q-1\in S\lor(n\text{ odd}\land S=E_n).
\]
The positive manuscript can use the same predicates and establish
the converse. Arity two is deliberately excluded.

\paragraph{Reusable logical lemmas.}
The one-point identification-or-strong dichotomy;
Lemma~\ref{lem:exactness} for unsigned cochains;
the parity count in Lemma~\ref{lem:zero-two};
and, crucially, the two-coordinate reflection
Lemma~\ref{lem:reflection}. The latter is common to both parity
branches and should be formalised only once.

\paragraph{Individual mathematical blocks.}
The two explicit arity-four witnesses (one implication chain and
one seventeen-row table); the arity-five defect support calculation
(including its sixteen-row incidence table); the one-point cone;
the even and rooted odd padding type maps;
the uniform defect-graph witness in arities $1\pmod4$;
and the disjoint-transversal counting proof in arities $3\pmod4$.
The closing induction should be formalised only after these pieces.

\paragraph{Genuine external mathematical dependencies.}
\begin{itemize}
\item \emph{No non-elementary external theorem is required by the
negative implication itself.} In particular, simplex cochain
exactness and the alternating-boundary identity are proved here.
\item The finite $n$-uniform Ramsey theorem is used only to explain
the equivalence between unboundedness and the presence of the central
local type. It can be omitted entirely from the Lean theorem by
taking $T_{n,q}\in S$ as an explicit hypothesis.
\item The probabilistic branch uses independence, finite counting,
the union bound and the elementary Bernoulli inequality. The
explicit choice of $t$ above avoids limits, exponential estimates
and other analytic prerequisites.
\item The positive AP/SAP constructions, the complete ternary census,
and any structural classification of homogeneous structures are
\emph{not} inputs to this proof.
\end{itemize}

\paragraph{Audit status and remaining checks.}
Independent finite enumeration confirms membership of both sides in
the intended classes for all three small seed diagrams. It verifies
all seventeen rows and thirty-four implication paths in the second
arity-four seed, all sixteen parity-support cases and five final
contradictions in the arity-five seed, and the uniform
$1\pmod4$ side diagrams for $n=5,9,13,17,21$.
Random-input checks confirm both local suspension type formulas at
several small arities. These finite checks complement the written
mathematical proofs, but do \emph{not} verify the infinite-family
lemmas or the counting argument mechanically. The full text remains
unaudited by a separate human referee and is not a Lean proof.

\begin{thebibliography}{9}
\bibitem{catalogue}
\emph{Local Types and Unbounded Amalgamation Classes for Constrained
Hypertournaments}, companion working classification manuscript,
2026 (nomenclature and context only; not a proof premise).
\bibitem{source-negative}
M.~Kone\v cn\'y,
\emph{Amalgamation of locally constrained hypertournaments:
a complete proof of the negative classification}, original
unpublished working note, 5 November 2026. Source text audited and
refactored in the present AI-written note.
\end{thebibliography}
\end{document}